% ============================================================================
% Chapter 13: Connecting to the Internet on Windows
% ============================================================================

\chapter{Connecting to the Internet}
\label{ch:internet-win}

\begin{objectives}
  \item Understand what an internet connection is
  \item Connect a Windows computer to a Wi-Fi network
  \item Check the status of your internet connection
  \item Troubleshoot basic connection problems
\end{objectives}

\bytetip[tip]{The internet connects your computer to the whole world. Let's learn how to get online --- and what to do if the connection drops!}

\section{Types of Internet Connections}
\label{sec:connection-types}
\index{internet!connection types}

\begin{theorybox}[Internet Connection Methods]
\begin{itemize}[leftmargin=*, label={\textcolor{ModuleBlue}{\faAngleRight}}, itemsep=3pt]
  \item \textbf{Wi-Fi}\index{Wi-Fi} --- Wireless connection via a router (most common at home and school)
  \item \textbf{Ethernet}\index{Ethernet} --- Wired connection using a cable (faster and more reliable)
  \item \textbf{Mobile Data}\index{mobile data} --- Uses your phone's cellular network (3G/4G/5G)
  \item \textbf{Broadband}\index{broadband} --- High-speed connection via phone line or fibre optic
\end{itemize}
\end{theorybox}

\section{Connecting to Wi-Fi}
\label{sec:connect-wifi}

To connect to Wi-Fi on Windows:
\begin{enumerate}[label=\protect\stepcircle{\arabic*}, leftmargin=2cm, itemsep=4pt]
  \item Click the \textbf{Wi-Fi icon} \faWifi\ in the System Tray (bottom-right).
  \item A list of available networks will appear.
  \item Click on your school or home network name.
  \item Enter the \textbf{password} if required and click Connect.
  \item A ``Connected'' message confirms you're online!
\end{enumerate}

\begin{notebox}
If you see a \faLock\ padlock icon next to a network name, it means the network is \textbf{password-protected}. You need the correct password to join. Open networks (without a padlock) don't require a password, but they are less secure.
\end{notebox}

\bytetip[warn]{Be careful with \textbf{public Wi-Fi} (in cafes, airports, malls). These networks are not secure --- someone could spy on your activity. Never log into your bank account or enter passwords on public Wi-Fi!}

\begin{funfact}
The word ``Wi-Fi'' doesn't actually stand for anything! Many people think it means ``Wireless Fidelity,'' but the Wi-Fi Alliance has confirmed that it's just a catchy brand name they invented. It sounds tech-y, and that was enough!
\end{funfact}

\begin{reviewqs}
  \item Name three types of internet connections.
  \item What steps do you follow to connect to Wi-Fi on Windows?
  \item What does a padlock icon next to a Wi-Fi network mean?
  \item Why should you be careful with public Wi-Fi?
  \item What is the difference between Wi-Fi and Ethernet?
\end{reviewqs}

\begin{challenge}
Check the Wi-Fi connection on your school computer. What is the network name? Is it password-protected? Open a web browser and go to \texttt{speedtest.net} to test your internet speed. Write down the \textbf{download speed} and \textbf{upload speed}.\end{challenge}
